\chapter{SPARC Rotation Curves: A Coefficient-Level Consistency Check}
\label{app:empirical}

\section{What is being tested}
\label{sec:sparc-what}

The function
\begin{equation}
  \label{eq:appD-nu}
  \nu_{\sqrt{}}(g_{\mathrm{N}},g_s)
  =\frac{1}{1-\exp\bigl(-\sqrt{g_{\mathrm{N}}/g_s}\bigr)}
\end{equation}
is the published Radial Acceleration Relation interpolator of
McGaugh, Lelli \& Schombert (2016). Appendix~\ref{app:status}
records that a derivation of this (or any other) interpolator from
the mutual-information functional remains open (Problem~H). The
present appendix asks only whether the same function, with the
single scale fixed at the published value
\(g_s=g_\dagger=1.20\times10^{-10}\,\mathrm{m\,s^{-2}}\),
is compatible with the SPARC rotation-curve sample.

The predicted circular velocity is
\begin{equation}
  \label{eq:appD-V}
  V=V_{\mathrm{b}}\sqrt{\nu_{\sqrt{}}(g_{\mathrm{N}},g_s)},
\end{equation}
where \(V_{\mathrm{b}}\) is the Newtonian baryonic contribution
computed from the SPARC mass models. No additional galactic
parameters are introduced.

\section{SPARC comparison}
\label{sec:sparc-comparison}

The comparison uses the public SPARC rotmod files
(175 galaxies, 3391 radial points). Three models are scored by
\(\chi^2\) and AIC:

\begin{itemize}
  \item baryons only (\(\nu=1\));
  \item simple MOND with one free acceleration scale \(a_0\);
  \item \(\nu_{\sqrt{}}\) with \(g_s\) fixed at \(g_\dagger\)
        (zero free parameters) and, for reference, with \(g_s\) free.
\end{itemize}

\begin{table}[htbp]
\centering
\caption{Global \(\chi^2\) and AIC on 175 SPARC galaxies (3391 points).
Lower AIC is better.}
\label{tab:sparc_results}
\begin{tabular}{lrrrr}
\toprule
Model & Free params & \(\chi^2\) & AIC & Scale (\(\mathrm{m\,s^{-2}}\)) \\
\midrule
\(\nu_{\sqrt{}}\) (fixed \(g_s=g_\dagger\)) & 0 & \(1.695\times10^5\) & \(1.695\times10^5\) & \(1.20\times10^{-10}\) \\
\(\nu_{\sqrt{}}\) (free \(g_s\))            & 1 & \(1.695\times10^5\) & \(1.695\times10^5\) & \(1.20\times10^{-10}\) \\
Simple MOND (free \(a_0\))                 & 1 & \(1.691\times10^5\) & \(1.691\times10^5\) & \(1.16\times10^{-10}\) \\
Baryons only                               & 0 & \(1.960\times10^6\) & \(1.960\times10^6\) & --- \\
\bottomrule
\end{tabular}
\end{table}

The fixed-\(g_s\) and free-\(g_s\) versions of \(\nu_{\sqrt{}}\) are
statistically indistinguishable; the free fit returns
\(g_s=1.20\times10^{-10}\,\mathrm{m\,s^{-2}}\). Relative to the
best one-parameter MOND fit the \(\chi^2\) penalty is
approximately \(0.2\%\). Both modified-gravity models improve on
baryons-only by more than an order of magnitude in \(\chi^2\).

\section{Numerical implementation}
\label{sec:sparc-code}

\begin{lstlisting}[style=namcgPythonStyle,caption={Core evaluation functions.}]
import numpy as np

KMS2_PER_KPC_TO_SI = 1.0e6 / 3.085677581491367e19
G_DAGGER = 1.20e-10

def gn_si(vb, r):
    r = np.maximum(np.asarray(r, float), 1e-6)
    return (np.asarray(vb, float)**2 / r) * KMS2_PER_KPC_TO_SI

def nu_sqrt(gn, gs):
    gs = max(float(gs), 1e-18)
    x = np.clip(np.asarray(gn, float)/gs, 0.0, 1e12)
    s = np.sqrt(x)
    denom = 1.0 - np.exp(-s)
    nu = np.empty_like(denom)
    tiny = denom < 1e-12
    nu[tiny] = 1.0 / np.maximum(s[tiny], 1e-30)
    nu[~tiny] = 1.0 / denom[~tiny]
    return nu

def v_namcg(r, vb, gs=G_DAGGER):
    return vb * np.sqrt(np.clip(nu_sqrt(gn_si(vb, r), gs), 1.0, 1e6))
\end{lstlisting}

\section{What the comparison is and is not}
\label{sec:sparc-status}

At low Newtonian acceleration,
\(\nu_{\sqrt{}}\) approaches the square-root map
\(g_{\mathrm{obs}}\approx\sqrt{g_{\mathrm{N}}g_\dagger}\).
On the SPARC disks this keeps \(V\) from falling as fast as
\(V_{\mathrm{b}}\). The numerical result shows that this
particular function, frozen at the published \(g_\dagger\),
is compatible with the sample at the same level as a
one-parameter MOND fit.

The comparison does not:
\begin{itemize}
  \item derive \(\nu_{\sqrt{}}\) from \(I(A:B)\) or from the
        modular spectrum (Problem~H remains open);
  \item fix \(g_\dagger\) from the condensate;
  \item constrain \(m_T=0\), the heat-kernel coefficients, or
        the conservation identity;
  \item rule out particle dark-matter halos on other scales.
\end{itemize}

It is a coefficient-level consistency check: the infrared
equation of the monograph can wear the published RAR
interpolator without contradiction on this data set. That is
the strongest statement the SPARC sample supports.